% Test-case specification TC-0028 (CTest test TC-0028-cases-varfuncs).
\documentclass[letterpaper]{article}

\usepackage[letterpaper,margin=1in]{geometry}
\emergencystretch=3em

% Compile from the repository root so this path resolves to
% the pinned WSP submodule.
\input{wsp/testing/test-case-library.tex}

\begin{document}

\def\TCID{TC-0028}
\def\TCTitle{%
  Call cases: varfuncs}
\def\TCPurpose{%
  Verify, against the installed runtime's own answers, \texttt{Abs},
  \texttt{Fix}, \texttt{Int}, \texttt{Sgn}, \texttt{Len}, \texttt{LenB},
  \texttt{LCase}, \texttt{UCase}, \texttt{Trim}, \texttt{LTrim} and
  \texttt{RTrim} on Variants (the cases where the reference raises error 13
  are not cases).}
\def\TCPriority{\TCPriorityHigh}

\def\TCRequirementRef{%
  REQ-0003, REQ-0004}
\def\TCDesignRef{%
  Not applicable.}
\def\TCFeatureRef{%
  \texttt{tests/\allowbreak{}vbrun\_\allowbreak{}helper\_\allowbreak{}tests.\allowbreak{}cpp}
  (family \texttt{cases}),
  \texttt{tests/\allowbreak{}cases/\allowbreak{}varfuncs.\allowbreak{}cases},
  \texttt{tests/\allowbreak{}cases/\allowbreak{}varfuncs.\allowbreak{}expected},
  \texttt{tests/\allowbreak{}cases/\allowbreak{}README.\allowbreak{}md}}
\def\TCTestTechnique{%
  Back-to-back comparison with the reference runtime's recorded results
  (\texttt{.\allowbreak{}expected}); argument kinds as partitions.}

\def\TCPreconditions{%
  The library and the test program are built
  (\texttt{scripts/\allowbreak{}Invoke-Build.\allowbreak{}ps1}) and the test
  is registered with CTest.}
\def\TCEnvironment{%
  Windows with the toolchain of
  \texttt{scripts/\allowbreak{}Invoke-Build.\allowbreak{}ps1}; the local
  gate builds and runs x64 Debug, CI builds and runs the other
  architectures. The library is built from this repository in
  \texttt{out/\allowbreak{}build}.}
\def\TCAssumptions{%
  Not applicable.}

\def\TCInputData{%
  The 164 call cases of
  \texttt{tests/\allowbreak{}cases/\allowbreak{}varfuncs.\allowbreak{}cases},
  one call of an export with its arguments per line (Strings, Longs,
  Variants, output Variants).}
\def\TCInitialState{%
  No state from earlier runs is needed; the test makes what it needs.}

\def\TCProcedure{%
  \begin{enumerate}
    \item CTest runs \texttt{vbrun\_\allowbreak{}helper\_\allowbreak{}tests
      \textless{}wfcvbrun.\allowbreak{}dll\textgreater{} cases
      varfuncs.\allowbreak{}cases varfuncs.\allowbreak{}expected}.
    \item Load the DLL and read the next case line (a \texttt{\#fixture} line
      rebuilds the scratch folder first).
    \item Call the named export with the arguments of the line and format the
      result as the README describes.
    \item Compare the formatted result with the same line of
      \texttt{varfuncs.\allowbreak{}expected}.
    \item CTest records the exit code and, on failure, the program's output
      (\texttt{--output-on-failure}).
  \end{enumerate}}
\def\TCExpectedResult{%
  All 164 results equal the recorded results; the program prints
  \texttt{cases ok: N checks} and exits 0. A difference prints the case, the
  expected and the actual text, and the exit code is 1.}
\def\TCPostConditions{%
  No file or registry state remains that later tests depend on.}

\TCTemplate

\end{document}
